\documentclass[a4paper, 12pt]{article}

%%% SST LAB PROTOCOLL PREAMBLE
%%% 2019
%%%%%%%%%%%%%%%%%%%%%%%%%%%%%%%


%%% PACKAGES
%%%%%%%%%%%%%%%%%%%%%%%%%%%

\usepackage[ngerman]{babel}

\usepackage[utf8]{inputenc}
\usepackage{amsmath}
\usepackage{pgfplots}
\usepackage{tikz}
\usepackage[many]{tcolorbox}
\usepackage{graphicx}
\graphicspath{ {./graphics/} }
\usepackage{pdfpages}
\usepackage{dashrule}
\usepackage{float}
\usepackage{siunitx}
\usepackage{trfsigns}
\usepackage{booktabs}
\usepackage[european]{circuitikz}
\usepackage{tcolorbox}

%%% DOCUMENT GEOMETRY
%%%%%%%%%%%%%%%%%%%%%%%%%%%

\usepackage{geometry}
\geometry{
 a4paper,
 total={0.6180339887498948\paperwidth,0.6180339887498948\paperheight},
 top = 0.1458980337503154\paperheight,
 bottom = 0.1458980337503154\paperheight
 }
\setlength{\jot}{0.013155617496424828\paperheight}
\linespread{1.1458980337503154}

\setlength{\parskip}{0.013155617496424828\paperheight} % paragraph spacing


%%% COLORS
%%%%%%%%%%%%%%%%%%%%%%%%%%%

\definecolor{red1}{HTML}{f38181}
\definecolor{yellow1}{HTML}{fce38a}
\definecolor{green1}{HTML}{95e1d3}
\definecolor{blue1}{HTML}{66bfbf}
\definecolor{hsblue}{HTML}{00b1db}
\definecolor{hsgrey}{HTML}{afafaf}

%%% CONSTANTS
%%%%%%%%%%%%%%%%%%%%%%%%%%%
\newlength{\smallvert}
\setlength{\smallvert}{0.0131556\paperheight}


%%% COMMANDS
%%%%%%%%%%%%%%%%%%%%%%%%%%%

% differential d
\newcommand*\dif{\mathop{}\!\mathrm{d}}

% horizontal line
\newcommand{\holine}[1]{
  	\begin{center}
	  	\noindent{\color{hsgrey}\hdashrule[0ex]{#1}{1pt}{3mm}}\\%[0.0131556\paperheight]
  	\end{center}
}

% mini section
\newcommand{\minisec}[1]{ \noindent\underline{\textit {#1} } \\}

% quick function plot
\newcommand{\plotfun}[3]{
  \vspace{0.021286\paperheight}
  \begin{center}
    \begin{tikzpicture}
      \begin{axis}[
        axis x line=center,
        axis y line=center,
        ]
        \addplot[draw=red1][domain=#2:#3]{#1};
      \end{axis}
    \end{tikzpicture}
  \end{center}
}

% box for notes
\newcommand{\notebox}[1]{

\tcbset{colback=white,colframe=green1!100!black,title=Note!,width=0.618\paperwidth,arc=0pt}

 \begin{center}
  \begin{tcolorbox}[]
   #1 
  \end{tcolorbox}
 
 \end{center} 
 
}

% box for equation
\newcommand{\eqbox}[2]{
	
	\tcbset{colback=white,colframe=green1!100!black,title=,width=#2,arc=0pt}
	
	\begin{center}
		\begin{tcolorbox}[ams align*]
				#1
		\end{tcolorbox}
		
	\end{center} 
	
}
% END OF PREAMBLE

%%%%%%%%%%%%%%%%%%%%%%%%%%%%%%%%%%%%%

\begin{document}

%%%%%%%%%%%%%%%%%%%%%%%%%%%%%%%%%%%%%
  \includepdf{./titlepage/titlepage.pdf}
  \clearpage
  \setcounter{page}{1}
%%%%%%%%%%%%%%%%%%%%%%%%%%%%%%%%%%%%%
\section{Prinzip der Frequenzmessung}

Durch Festlegen einer \emph{Torzeit} (Gate Time) $T_{\textrm{Gate}}$ und Zählen der
Flanken des Eingangssignals ($Z$) während dieser Zeit, kann die Frequenz des Signals
berechnet werden.
\[f = \frac{Z}{T_{\textrm{Gate}}}\]

\section{Umsetzung}
Das zu messende Eingangssignal liegt an P2.2 an, welcher zur Capture-Einheit 0
des Timers A gehört. Dort muss dann im TACCTL-Register über das CCIS-Flag der
Eingang CCI0B (01b) ausgewählt werden. Die Torzeit wird dann mit dem Timer B
generiert. Die Berechnung wird im Interrupt (TAIFG) des Timer B durchgeführt.

\section{Erzeugung der Torzeit}
Ist der Timer im continuous-Mode, lässt sich bei einer Taktfrequenz von 16 MHz
und dem maximalen Teilerwert von 8 eine maximale Periodendauer von
\[T_{max} = \frac{2^{16}-1}{16/8 \,\si{\mega\hertz}} = 32.7675 \, \si{\milli\second}\]

einstellen. Um eine Torzeit von $1 \, \si{\second}$ zu erreichen, muss man nun
die Anzahl der Timerüberläufe (Interrupt TAIFG) zählen, bis diese den nötigen Wert $k$ erreicht haben.
Dieser ist für $T_{max}$
\[k = \frac{1 \, \si{\second}}{T_{max}} = 30.518\]

Da $k$ in diesem Fall kein ganzzahliger Wert ist, man aber nur ganze Überläufe
zählen kann, entstehen Ungenauigkeiten. Um das zu verhindern, kann man den Timer
auch in den up-Mode setzen und zu einem Wert zählen lassen, der bei der
gewünschten Zeit eine ganzzahlige Anzahl an Überläufen/Interrupts erzeugt,
beispielsweise 40 Überläufe, woraus sich eine Periodendauer von
\[T_{up} = \frac{1 \, \si{\second}}{40} = 25 \, \si{\milli\second}\]
ergibt.

Um diese Zeit zu erzeugen, braucht der Timer im up-Mode den Comparewert
\[\textrm{CCR0} = 25 \, \si{\milli\second} \cdot 2 \, \si{\mega\hertz} = 50000\]

\section{Genauigkeit}
Die sich ergebende Messabweichung ist
\[F = \pm \frac{1}{2} \cdot \frac{1}{T_{\textrm{Gate}}}\]

\section{Portinitialisierung}
P2.2 muss für die Input-Richtung konfiguriert werden (0).
\begin{lstlisting}
P2DIR &= ~(1<<2);
\end{lstlisting}

\noindent Außerdem muss die Zweitfunktion ausgewählt werden (1).
\begin{lstlisting}
P2SEL |= (1<<2);
\end{lstlisting}


\end{document}

% Local Variables:
% TeX-engine: luatex
% End:
